%!TEX root = ../template.tex
%%%%%%%%%%%%%%%%%%%%%%%%%%%%%%%%%%%%%%%%%%%%%%%%%%%%%%%%%%%%%%%%%%%%
%% abstract-pt.tex
%% NOVA thesis document file
%%
%% Abstract in Portuguese
%%%%%%%%%%%%%%%%%%%%%%%%%%%%%%%%%%%%%%%%%%%%%%%%%%%%%%%%%%%%%%%%%%%%

\typeout{NT FILE abstract-pt.tex}%

O software de open-source e os rádios definidos por software de baixo custo permitem implementar uma rede 5G privada em modo Standalone sem o orçamento de um operador comercial. Implementá-la, porém, é apenas metade do problema, pois uma rede funcional pode ainda ficar muito aquém do que o equipamento permite, e calibrá-la às condições em que opera é um outro desafio, ainda mais difícil sem uma referência.

Esta dissertação constrói uma rede 5G Standalone, juntamente com uma rede 4G, para referência, partilhando o mesmo núcleo Open5GS. A rede 5G utiliza o srsRAN Project e uma USRP B210, e a rede 4G o srsRAN 4G e uma bladeRF xA4. As duas foram comparadas a três distâncias do rádio, ajustando a rede 5G ao longo de três rondas.

Face a esta referência, o uplink 5G sem ajustes era mais de 20 vezes mais lento a 10\,m. Mudar o padrão TDD de 3:1 para 2:2 igualou o throughput ao da rede 4G, e reduzir o ganho do recetor eliminou a saturação do equipamento a curta distância, atingindo 28,5\,Mbps a 2\,m, o dobro da 4G, mas à custa da cobertura a 10\,m. O downlink 5G atingiu 50,6\,Mbps, mais do dobro do 4G. O tempo de ida e volta médio, de 25 a 28\,ms sem carga, subiu para segundos sob tráfego TCP, devido a uma queue RLC sobredimensionada.

Estes resultados mostram que uma rede 5G Standalone open-source pode igualar ou superar uma rede 4G construída com o mesmo equipamento, mas só depois de ajustada às condições que pretende-se operar, e que nenhuma configuração fixa serve todos as localizações. A plataforma fica instalada e documentada para trabalho futuro.

% Palavras-chave do resumo em Português
\keywords{
  5G Standalone \and
  Open-source \and
  5G Testbed \and
  Software-Defined Radio \and
  Avaliação de desempenho
}
